\chapter{محیط (The Environment)}

همان‌طور که پیش‌تر اشاره شد، پوسته (\en{Shell}) در طول هر نشست (\en{Session})، مجموعه‌ای از داده‌های حیاتی را در ساختاری تحت عنوان «محیط (\en{Environment})» نگهداری می‌کند. این محیط دربرگیرنده متغیرهایی است که برنامه‌ها به‌منظور پیکربندی و تنظیم رفتار خود به آن‌ها رجوع می‌کنند. با درک عمیق مکانیزم‌های حاکم بر این محیط، قادر خواهیم بود تجربه کاربری خود را در محیط پوسته به‌طور مؤثری شخصی‌سازی کنیم.

در این فصل، برای تعامل با محیط و مدیریت آن، دستورات کلیدی زیر را بررسی خواهیم کرد:

\begin{itemize}
  \item \cmd{printenv}: جهت نمایش بخشی یا تمامی متغیرهای محیطی (\en{Environment Variables}).
  \item \cmd{set}: برای نمایش، تنظیم یا غیرفعال‌سازی گزینه‌های پوسته (\en{shell options}) و همچنین مدیریت متغیرهای محلی و پارامترهای موقعیتی.
  \item \cmd{export}: به‌منظور صدور یک متغیر از پوسته به محیط، جهت دسترسی فرآیندهای فرزند (\en{Child Processes}) به آن.
  \item \cmd{alias}: جهت تعریف نام‌های مستعار برای دستورات یا رشته‌ای از فرمان‌های ترکیبی.
  \item \cmd{source}: برای بارگذاری و اجرای مستقیم دستورات موجود در یک فایل درون نشست فعلی پوسته.
\end{itemize}

\section{انواع داده‌های محیطی}

پوسته دو نوع اصلی از داده‌ها را در محیط خود نگهداری می‌کند که اگرچه در نگاه اول در \en{Bash} رفتاری مشابه دارند، اما از دیدگاه قلمرو دسترسی (\en{Scope}) تفاوت‌های بنیادی با یکدیگر دارند:

\begin{itemize}
  \item \textbf{متغیرهای پوسته (\en{Shell Variables}):} این متغیرها منحصراً توسط نمونه فعلیِ پوسته ایجاد و مدیریت می‌شوند و تنها در دامنه همان نشست (\en{Session}) جاری معتبر هستند.
  \item \textbf{متغیرهای محیطی (\en{Environment Variables}):} این متغیرها در پوستهٔ والد (\en{Parent Shell}) تعریف می‌شوند و با استفاده از دستور \cmd{export} به محیط صادر می‌گردند تا به‌طور خودکار به تمامی فرآیندهای فرزند (\en{Child Processes}) که از آن پوسته مشتق می‌شوند، به ارث برسند.
\end{itemize}

علاوه بر متغیرهای یادشده، پوسته اجزای دیگری را نیز جهت تسهیل فرآیند اتوماسیون و بهینه‌سازی عملیات در حافظه ذخیره می‌کند:

\begin{itemize}
  \item \textbf{نام‌های مستعار (\en{Aliases}):} میان‌برهایی برای دستورات طولانی که در فصل پنجم بررسی شدند.
  \item \textbf{توابع پوسته (\en{Shell Functions}):} که در بخش چهارم به‌صورت تخصصی به آن‌ها خواهیم پرداخت.
\end{itemize}

این ساختار داده‌ای پویا به پوسته و سایر برنامه‌های وابسته اجازه می‌دهد تا پیکربندی و رفتار خود را بر اساس تنظیمات سیستمی و اطلاعات لحظه‌ای، به‌صورت یکپارچه و هماهنگ تنظیم کنند.

\section{پایش و بررسی محیط (Environment)}

برای بازبینی محتویات محیط، دو ابزار کلیدی در اختیار دارید: دستور داخلی \cmd{set} در پوسته \en{Bash} و برنامه مستقل \cmd{printenv}. این دو دستور از نظر دامنه نمایش تفاوت‌های بنیادینی با یکدیگر دارند:

\begin{itemize}
  \item \cmd{set}: این دستور جامع‌تر عمل کرده و تمامی متغیرهای پوسته (\en{Shell Variables})، متغیرهای محیطی (\en{Environment Variables}) و همچنین توابع (\en{Functions}) تعریف‌شده را نمایش می‌دهد.
  \item \cmd{printenv}: این دستور منحصراً بر نمایش متغیرهای محیطی تمرکز دارد.
\end{itemize}

از آنجایی که حجم داده‌های موجود در محیط معمولاً بسیار زیاد است، توصیه می‌شود خروجی را جهت پیمایش آسان‌تر به دستور \cmd{less} پایپ (\en{Pipe}) کنید:

\begin{latin}
\begin{lstlisting}
[me@linuxbox ~]$ printenv | less
Doing so, we should get something
that looks like this:
USER=me
PAGER=less
LSCOLORS=6xfxcxdxbxegedabagacad
XDG_CONFIG_DIRS=/etc/xdg/xdg-ubuntu:/usr/share/upstart/xdg:/etc/xdg
PATH=/home/me/bin:/usr/local/sbin:/usr/local/bin:/usr/sbin:/usr/bin:/sbin:/bin:/usr/games:/usr/local/games
DESKTOP_SESSION=ubuntu
QT_IM_MODULE=ibus
QT_QPA_PLATFORMTHEME=appmenu-qt5
JOB=dbus
PWD=/home/me
GNOME_KEYRING_PID=1850
LANG=en_US.UTF-8
GDM_LANG=en_US
MANDATORY_PATH=/usr/share/gconf/ubuntu.mandatory.path
MASTER_HOST=linuxbox
IM_CONFIG_PHASE=1
COMPIZ_CONFIG_PROFILE=ubuntu
GDMSESSION=ubuntu
SESSIONTYPE=gnome-session
XDG_SEAT=seat0
HOME=/home/me
SHLVL=2
LANGUAGE=en_US
GNOME_DESKTOP_SESSION_ID=this-is-deprecated
LESS=-R
LOGNAME=me
COMPIZ_BIN_PATH=/usr/bin/
LC_CTYPE=en_US.UTF-8
XDG_DATA_DIRS=/usr/share/ubuntu:/usr/share/gnome:/usr/local/share://usr/share/
QT4_IM_MODULE=xim
DBUS_SESSION_BUS_ADDRESS=unix:abstract=/tmp/dbus-IwaesmWaT0
LESSOPEN=| /usr/bin/lesspipe %s
INSTANCE=
\end{lstlisting}
\end{latin}

این خروجی، مجموعه‌ای از متغیرها و مقادیر متناظر آن‌ها را نمایش می‌دهد که در آن هر سطر از فرمت استاندارد \cmd{NAME=value} پیروی می‌کند. برنامه‌ها از این متغیرها برای استخراج تنظیمات، شناسایی مسیرهای اجرایی و دریافت اطلاعات کلیدی سیستم استفاده می‌کنند.

جهت استخراج مقدار یک متغیر محیطیِ معین، دستور \cmd{printenv} را به‌همراه نام آن متغیر اجرا می‌کنیم. این شیوه، سریع‌ترین راه برای دسترسی مستقیم به داده‌های پیکربندی‌شده در محیط عملیاتی است.

به عنوان نمونه، برای مشاهده نام کاربری جاری:

\begin{latin}
\begin{lstlisting}
[me@linuxbox ~]$ printenv USER
me
\end{lstlisting}
\end{latin}

دستور داخلی \cmd{set} در پوسته، چنانچه بدون هیچ آرگومانی اجرا شود، فهرستی جامع شامل تمامی متغیرهای پوسته، متغیرهای محیطی و همچنین توابع (\en{Functions}) تعریف‌شده را به‌ترتیب حروف الفبا نمایش می‌دهد. با توجه به حجم بسیار زیاد این خروجی، استفاده از یک پایپ‌لاین (\en{Pipe}) به دستور \cmd{less} برای مدیریت پیمایش توصیه می‌شود:

\begin{latin}
\begin{lstlisting}
[me@linuxbox ~]$ set | less
\end{lstlisting}
\end{latin}

علاوه بر این، برای استخراج مقدار یک متغیر خاص، می‌توان از دستور \cmd{echo} بهره برد. در این روش، باید نام متغیر را با پیشوند \cmd{\$} همراه کنید تا پوسته آن را به عنوان یک مرجع متغیر (\en{Variable Reference}) شناسایی و مقدارش را جایگزین کند:

\begin{latin}
\begin{lstlisting}
[me@linuxbox ~]$ echo $HOME
/home/me
\end{lstlisting}
\end{latin}

نام‌های مستعار یا \en{Aliases} از جمله اجزای محیطی هستند که در خروجی دستورات \cmd{set} یا \cmd{printenv} ظاهر نمی‌شوند. برای بازبینی فهرست آن‌ها، کافی است دستور \cmd{alias} را بدون پارامتر ورودی اجرا نمایید:

\begin{latin}
\begin{lstlisting}
[me@linuxbox ~]$ alias
alias l='ls -d .* --color=tty'
alias ll='ls -l --color=tty'
alias ls='ls --color=tty'
alias vi='vim'
alias which='alias | /usr/bin/whi\ch --tty-only --read-alias --show-dot --show-tilde'
\end{lstlisting}
\end{latin}

خروجی فوق، تمامی نام‌های مستعار تعریف‌شده در نشست جاری را نمایش می‌دهد؛ ابزارهایی که با جایگزینی دستورات پیچیده و طولانی با عناوین کوتاه، بهره‌وری در محیط خط فرمان را افزایش می‌دهند.

\section{متغیرهای مهم محیطی}

محیط عملیاتی پوسته دربرگیرنده متغیرهای متعددی است که وضعیت و مشخصات نشست جاری را ثبت می‌کنند. اگرچه جزئیات این متغیرها در توزیع‌های مختلف لینوکس تفاوت‌هایی دارد، اما مجموعه‌ای از آن‌ها به‌عنوان استانداردهای رایج در اکثر محیط‌ها حضور دارند. جدول زیر به تشریح برخی از کاربردی‌ترینِ این متغیرها می‌پردازد:

\begin{table}[htbp]
\centering
\caption{متغیرهای محیطی متداول}
\label{tab:common-env-vars}
\begin{tabularx}{\textwidth}{l X}
\toprule
\textbf{متغیر} & \textbf{مفهوم و محتوا} \\
\midrule
\cmd{DISPLAY} & شناسهٔ نمایشگر فعال در محیط گرافیکی. این متغیر عمدتاً برای برنامه‌های مبتنی بر \en{X11} کاربرد دارد.
\begin{itemize}
  \item \textbf{در محیط‌های \en{X11}:} مقدار \cmd{:0} معمولاً به اولین نمایشگر متصل به سرور \en{X11} اشاره دارد.
  \item \textbf{در محیط‌های \en{Wayland}:} این متغیر توسط برنامه‌های بومی \en{Wayland} استفاده نمی‌شود و ممکن است مقدار آن خالی باشد. با این حال، اگر مقدار \cmd{:0} یا \cmd{:1} را مشاهده کنید، نشان‌دهندهٔ وجود لایهٔ سازگاری \en{XWayland} است که برای اجرای برنامه‌های قدیمی \en{X11} در نظر گرفته شده است. برنامه‌های بومی \en{Wayland} برای اتصال به سرور نمایشگر از متغیر \cmd{WAYLAND\_DISPLAY} (که معمولاً مقدار \cmd{wayland-0} یا \cmd{wayland-1} دارد) استفاده می‌کنند.
\end{itemize} \\
\addlinespace
\cmd{EDITOR} & مسیر یا نام ویرایشگر متن پیش‌فرض سیستم جهت ویرایش فایل‌های متنی. \\
\addlinespace
\cmd{SHELL} & مسیر اجراییِ پوسته پیش‌فرض کاربر (مانند \file{/bin/bash} یا \file{/bin/zsh}). \\
\addlinespace
\cmd{HOME} & مسیر مطلق دایرکتوری خانگی (\en{Home Directory}) کاربر جاری. \\
\addlinespace
\cmd{LANG} & تنظیمات بومی‌سازی (\en{Locale}) شامل زبان، منطقه جغرافیایی و یونیکد (مانند \cmd{en\_US.UTF-8}). \\
\addlinespace
\cmd{OLDPWD} & مسیر آخرین دایرکتوری که کاربر پیش از جابه‌جایی فعلی در آن حضور داشته است. این متغیر توسط دستور \cmd{cd -} استفاده می‌شود. \\
\addlinespace
\cmd{PAGER} & ابزار پیش‌فرض جهت صفحه‌بندی و مرور خروجی‌های طولانی (مانند \cmd{less}). \\
\addlinespace
\cmd{PATH} & حیاتی‌ترین متغیر محیطی؛ فهرستی از دایرکتوری‌ها که پوسته برای یافتن فایل‌های اجرایی، آن‌ها را به‌ترتیب جستجو می‌کند. \\
\addlinespace
\cmd{PS1} & تعریف‌کننده ساختار و ظاهر اعلان خط فرمان (\en{Prompt}). این متغیر انعطاف بالایی در شخصی‌سازی دارد. \\
\addlinespace
\cmd{PWD} & مسیر مطلق دایرکتوری کاری فعلی (\en{Present Working Directory}). \\
\addlinespace
\cmd{TERM} & این متغیر، شناسه‌ای متنی است که برنامه‌ها از آن برای تشخیص قابلیت‌های ترمینال (مانند تعداد رنگ‌های قابل‌نمایش، امکان جابه‌جایی مکان‌نما، پشتیبانی از ویژگی‌های نمایشی مانند پررنگ، ایتالیک و زیرخط) استفاده می‌کنند. مقدار رایج در توزیع‌های مدرن لینوکس، \cmd{xterm-256color} است که پشتیبانی از ۲۵۶ رنگ را فعال می‌کند. \\
\addlinespace
\cmd{TZ} & منطقه زمانی (\en{Time Zone})؛ جهت تبدیل ساعت هماهنگ جهانی (\en{UTC}) به زمان محلی. \\
\addlinespace
\cmd{USER} & نام کاربریِ صاحب نشست جاری. \\
\bottomrule
\end{tabularx}
\end{table}

شایان ذکر است که عدم مشاهده برخی از این متغیرها در محیط شما به معنای نقص سیستم نیست؛ چرا که پیکربندی این مقادیر مستقیماً به نوع توزیع (\en{Distribution}) و تنظیمات اختصاصی حساب کاربری شما بستگی دارد.

\section{فرآیند ایجاد و شکل‌گیری محیط}

هنگام ورود به سیستم، مفسر \en{Bash} فراخوانی شده و فرآیند پیکربندی محیط کاربر را از طریق بارگذاری مجموعه‌ای از فایل‌های آغازین (\en{Startup Files}) آغاز می‌کند. این فرآیند به‌صورت سلسله‌مراتبی طی دو مرحله اجرا می‌شود: ابتدا، فایل‌های پیکربندی سراسری (\en{Global}) که بر تمامی کاربران سیستم حاکم است فراخوانی می‌شوند و سپس، فایل‌های اختصاصی در دایرکتوری خانگی کاربر پردازش می‌گردند تا تنظیمات شخصی‌سازی شده اعمال شود.

ترتیب و نحوه اجرای این فایل‌ها مستقیماً به نوع نشست پوسته بستگی دارد که به دو دسته اصلی تقسیم می‌شوند:

\begin{itemize}
  \item \textbf{نشست پوسته‌ی ورود (\en{Login Shell Session}):} این نشست زمانی شکل می‌گیرد که کاربر از طریق فرآیند احراز هویت (\en{Login})، به‌طور مستقیم به سیستم متصل شود. ورود به کنسول‌های مجازی (\en{Virtual Consoles}) یا برقراری اتصال از راه دور تحت پروتکل \en{SSH}، بارزترین نمونه‌های این نوع نشست هستند. در یک نشست ورود، \en{Bash} ابتدا فایل پیکربندی سراسری \file{/etc/profile} را در صورت وجود، بارگذاری و اجرا می‌کند. سپس، به‌ترتیب به جستجوی فایل‌های \file{\textasciitilde/.bash\_profile}، \file{\textasciitilde/.bash\_login} و \file{\textasciitilde/.profile} در دایرکتوری خانگی کاربر پرداخته و اولین فایلی که موجود و قابل خواندن باشد را اجرا می‌نماید.

  \item \textbf{نشست پوسته‌ی غیرورود (\en{Non-login Shell Session}):} به نشستی گفته می‌شود که بدون طی کردن مجدد فرآیند احراز هویت (ورود با نام کاربری و رمز عبور) در یک نشست فعال راه‌اندازی می‌شود. این نوع نشست، برخلاف نشست ورود (\en{Login Shell})، فایل‌های پروفایل مانند \file{/etc/profile}، \file{\textasciitilde/.bash\_login}، \file{\textasciitilde/.bash\_profile} یا \file{\textasciitilde/.profile} را اجرا نمی‌کند. در عوض، یک پوسته‌ی غیرورودِ تعاملی (\en{Interactive Non-login Shell}) برای راه‌اندازی، ابتدا فایل پیکربندی سراسری \file{/etc/bash.bashrc} را در صورت وجود، بارگذاری و اجرا کرده و سپس فایل اختصاصی کاربر \file{\textasciitilde/.bashrc} را فراخوانی می‌نماید. نمونه‌های رایج این نوع نشست عبارتند از: باز کردن یک شبیه‌ساز ترمینال (\en{Terminal Emulator}) در بستر یک محیط گرافیکی فعال (مانند \en{GNOME} یا \en{KDE})، اجرای دستور \cmd{bash} درون یک پوسته‌ی موجود برای ایجاد یک زیرپوسته (\en{Subshell})، تغییر هویت با دستور \cmd{su} به‌صورت بدون خط تیره (\cmd{su username})، اجرای یک اسکریپت پوسته (مانند \cmd{bash script.sh} یا \cmd{./script.sh}) یا اجرای یک دستور از راه دور از طریق \en{SSH} (مانند \cmd{ssh user@host command}) که یک پوسته‌ی غیرتعاملی و غیرورود ایجاد می‌کند.
\end{itemize}

همان‌طور که اشاره شد، پوسته‌های ورود، فایل‌های پیکربندی آغازین خود را بر اساس اولویت و توالی مندرج در جدول زیر بارگذاری می‌کنند:

\begin{table}[htbp]
\centering
\caption{فایل‌های آغازین نشست‌های پوسته ورود}
\label{tab:login-shell-files}
\begin{tabularx}{\textwidth}{l X}
\toprule
\textbf{فایل} & \textbf{محتوا و عملکرد} \\
\midrule
\file{/etc/profile} & اسکریپت پیکربندی سراسری سیستم که در ابتدای هر نشست ورود، توسط پوسته اجرا می‌شود. این فایل، متغیرهای محیطی پایه و تنظیمات مشترک تمام کاربران را تعیین می‌کند. \\
\addlinespace
\file{\textasciitilde/.bash\_profile} & نخستین فایل راه‌انداز اختصاصی کاربر که در دایرکتوری خانگی او قرار دارد. در صورت وجود، پوسته آن را پس از فایل سراسری اجرا کرده و امکان شخصی‌سازی محیط و بازنویسی تنظیمات سراسری را فراهم می‌کند. \\
\addlinespace
\file{\textasciitilde/.bash\_login} & در صورت عدم وجود فایل \file{\textasciitilde/.bash\_profile}، پوسته \en{Bash} به‌صورت خودکار برای یافتن دستورات پیکربندی به سراغ این فایل می‌رود. \\
\addlinespace
\file{\textasciitilde/.profile} & آخرین گزینه در زنجیرهٔ جستجوی فایل‌های پروفایل کاربری. چنانچه هیچ‌یک از دو فایل پیشین \file{\textasciitilde/.bash\_profile} و \file{\textasciitilde/.bash\_login} موجود نباشند، پوسته این فایل را اجرا می‌کند. این فایل در توزیع‌های مبتنی بر دبیان (مانند اوبونتو) از اهمیت ویژه‌ای برخوردار است و معمولاً به‌عنوان فایل اصلی پیکربندی کاربر در نظر گرفته می‌شود. \\
\bottomrule
\end{tabularx}
\end{table}

نشست‌های پوسته غیرورود (\en{Non-login Shells})، فایل‌های پیکربندی مندرج در جدول زیر را فراخوانی می‌کنند:

\begin{table}[htbp]
\centering
\caption{فایل‌های آغازین نشست‌های پوسته غیرورود}
\label{tab:non-login-shell-files}
\begin{tabularx}{\textwidth}{l X}
\toprule
\textbf{فایل} & \textbf{محتوا و عملکرد} \\
\midrule
\file{/etc/bash.bashrc} & اسکریپت پیکربندی سراسری سیستم (\en{System-wide}) که تنظیمات محیطی مشترک را برای تمامی نشست‌های غیرورودی کاربران اعمال می‌کند. \\
\addlinespace
\file{\textasciitilde/.bashrc} & فایل پیکربندی اختصاصی کاربر که در دایرکتوری خانگی او قرار دارد. این فایل پس از فایل سراسری (در صورت وجود) بارگذاری شده و وظیفهٔ شخصی‌سازی محیط پوسته، تعریف نام‌های مستعار، توابع و بازنویسی تنظیمات سراسری را بر عهده دارد. \\
\bottomrule
\end{tabularx}
\end{table}

\begin{note}
\textbf{نکته:} در نسخهٔ خالص \en{Bash}، فایل \file{/etc/bash.bashrc} به‌صورت پیش‌فرض وجود ندارد و با فعال‌سازی گزینهٔ ویژه در زمان کامپایل در دسترس قرار می‌گیرد. با این حال، اکثر توزیع‌های مدرن لینوکس از آن پشتیبانی کرده و آن را به‌عنوان بخشی از پیکربندی استاندارد خود ارائه می‌دهند. در برخی توزیع‌ها، این فایل سراسری به جای \file{/etc/bash.bashrc} با نام \file{/etc/bashrc} یا \file{/etc/bash.bashrc.local} ذخیره می‌شود.
\end{note}

پوسته‌های غیرورودی علاوه بر بارگذاری فایل‌های پیکربندی مندرج در جدول فوق، متغیرهای محیطی (\en{Environment Variables}) را از فرآیند والد خود ــ که غالباً یک نشست ورودی است ــ به ارث می‌برند.

توصیه می‌شود جهت درک بهتر پیکربندی سیستم خود، این فایل‌های آغازین را بررسی نمایید. توجه داشته باشید که به دلیل شروع نام این فایل‌ها با نویسه نقطه (\cmd{.}) که نشان‌دهنده پنهان بودن (\en{Hidden}) آن‌هاست، هنگام استفاده از دستور \cmd{ls} باید از گزینه \cmd{-a} جهت نمایش آن‌ها استفاده کنید.

از دیدگاه یک کاربر عادی، فایل \file{\textasciitilde/.bashrc} کلیدی‌ترین فایل پیکربندی محسوب می‌شود؛ چرا که این فایل تقریباً در تمامی سناریوها فراخوانی می‌گردد. نشست‌های غیرورودی به‌طور پیش‌فرض آن را اجرا می‌کنند و ساختار اکثر فایل‌های آغازین در نشست‌های ورودی نیز به‌گونه‌ای است که این فایل را در حین اجرا فراخوانی (\en{Source}) می‌کنند.

\section{تحلیل ساختار و محتوای فایل راه‌انداز}

با بررسی محتوای یک فایل \file{.bash\_profile} استاندارد (استخراج شده از توزیع \en{CentOS 6})، ساختاری مشابه نمونه زیر را مشاهده خواهیم کرد:

\begin{latin}
\begin{lstlisting}
# .bash_profile

# Get the aliases and functions
if [ -f ~/.bashrc ]; then
    . ~/.bashrc
fi

# User specific environment and startup programs

PATH=$PATH:$HOME/bin
export PATH
\end{lstlisting}
\end{latin}

خطوطی که با نویسه \cmd{\#} آغاز شده‌اند، کامنت (توضیحات) هستند و توسط پوسته نادیده گرفته می‌شوند؛ این متون صرفاً جهت ارتقای خوانایی اسکریپت برای کاربر درج شده‌اند.

نخستین بخش حائز اهمیت در خط چهارم قرار دارد:

\begin{latin}
\begin{lstlisting}
if [ -f ~/.bashrc ]; then
    . ~/.bashrc
fi
\end{lstlisting}
\end{latin}

در این قطعه کد، \cmd{if} یک دستور مرکب (\en{Compound Command}) است که در بخش چهارم این کتاب و در مبحث اسکریپت‌نویسی پوسته به‌صورت تفصیلی بررسی خواهد شد. با این حال، منطق اجرایی آن در مقطع کنونی بدین شرح است: چنانچه فایل \file{\textasciitilde/.bashrc} در مسیر مربوطه موجود باشد، محتویات آن را خوانده و اجرا کن.

\begin{latin}
\begin{lstlisting}
If the file "~/.bashrc" exists, then
     read the "~/.bashrc" file.
\end{lstlisting}
\end{latin}

مشاهده می‌کنیم که این سازوکار، روشی است که یک «پوسته ورود» (\en{Login Shell}) از طریق آن محتویات فایل \file{.bashrc} را فراخوانی یا به‌اصطلاح فنی \en{Source} می‌کند.

بخش بعدی در فایل پیکربندی به متغیر محیطی \cmd{PATH} اختصاص دارد. آیا تاکنون تأمل کرده‌اید که پوسته چگونه محل دقیق دستوری مانند \cmd{ls} را شناسایی می‌کند؟ در واقع، هنگام وارد کردن دستور \cmd{ls} در خط فرمان، پوسته تمام سیستم فایل را برای یافتن فایل اجرایی جستجو نمی‌کند؛ بلکه صرفاً فهرستی از دایرکتوری‌های تعریف‌شده در متغیر \cmd{PATH} را بررسی می‌نماید تا برنامه مورد نظر را بیابد.

متغیر \cmd{PATH} اغلب (هرچند نه همیشه، و بسته به نوع توزیع لینوکس) در فایل پیکربندی \file{/etc/profile} و به روش زیر مقداردهی می‌شود:

\begin{latin}
\begin{lstlisting}
PATH=$PATH:$HOME/bin
\end{lstlisting}
\end{latin}

در این عبارت، متغیر \cmd{PATH} به‌گونه‌ای بازتعریف می‌شود که دایرکتوری \file{bin/} موجود در دایرکتوری خانگی کاربر (\cmd{\$HOME}) به انتهای فهرست مسیرهای جستجو اضافه شود. این فرآیند، نمونه‌ای از «بسط پارامتر» (\en{Parameter Expansion}) است که پیش‌تر در فصل هفتم به بررسی آن پرداختیم.

برای درک بهتر سازوکار این فرآیند، دستورات زیر را امتحان کنید:

\begin{latin}
\begin{lstlisting}
[me@linuxbox ~]$ foo="This is some "
[me@linuxbox ~]$ echo $foo
This is some
[me@linuxbox ~]$ foo=$foo"text."
[me@linuxbox ~]$ echo $foo
This is some text.
\end{lstlisting}
\end{latin}

همان‌طور که در مثال فوق مشاهده می‌شود، با بهره‌گیری از این تکنیک می‌توان به‌سادگی یک رشته را به انتهای مقدار یک متغیر الحاق (\en{Append}) کرد. این همان مکانیزمی است که در مقداردهی متغیر \cmd{PATH} به کار رفته است.

با افزودن عبارت \cmd{\$HOME/bin} به انتهای متغیر \cmd{PATH} در فایل‌های پیکربندی، دایرکتوری مذکور به فهرست مسیرهایی که پوسته هنگام جستجوی دستورات بررسی می‌کند، اضافه می‌شود. این یعنی سیستم به‌گونه‌ای پیکربندی شده است که به‌طور پیش‌فرض از برنامه‌های اختصاصی کاربر پشتیبانی کند. بنابراین، اگر دایرکتوری \file{bin} را در دایرکتوری خانگی خود ایجاد کرده و اسکریپت‌ها یا فایل‌های اجرایی خود را در آن قرار دهید، پوسته به‌طور خودکار آن‌ها را شناسایی کرده و بدون نیاز به وارد کردن مسیر کامل، قابل اجرا خواهند بود.

در نهایت، دستور \cmd{export} تضمین می‌کند که این تغییرات در تمامی زیرپوسته‌ها (\en{Subshells}) و فرآیندهای مشتق‌شده از پوستهٔ جاری نیز اعمال گردند و متغیر \cmd{PATH} اصلاح‌شده، به عنوان یک متغیر محیطی (\en{Environment Variable}) در دسترس آن‌ها قرار گیرد.

\begin{latin}
\begin{lstlisting}
export PATH
\end{lstlisting}
\end{latin}

\begin{note}
\textbf{نکته:} بسیاری از توزیع‌های لینوکس این پیکربندی را به‌صورت پیش‌فرض در اختیار کاربر قرار می‌دهند. برای مثال، توزیع‌های مبتنی بر دبیان (مانند اوبونتو)، در هنگام فرآیند ورود (\en{Login})، وجود دایرکتوری \file{\textasciitilde/bin} را بررسی کرده و در صورت یافتن آن، این مسیر را به‌صورت پویا (\en{Dynamic}) به متغیر \cmd{PATH} الحاق می‌کنند.
\end{note}

\section{وراثت محیطی در فرآیندهای فرزند}

نکته پایانی در بخش پیشین مستلزم تبیین دقیق‌تری است. به‌طور کلی، متغیرهای پوسته تنها در نمونه‌ی جاری (\en{Instance}) معتبر هستند و به‌صورت خودکار به فرآیندهای فرزندی که توسط آن پوسته ایجاد می‌شوند، منتقل نمی‌گردند. برای درک بهتر این مفهوم، سناریوی زیر را بررسی می‌کنیم.

ابتدا یک متغیر محلی در نشست جاری پوسته تعریف می‌کنیم:

\begin{latin}
\begin{lstlisting}
[me@linuxbox ~]$ foo="bar"
\end{lstlisting}
\end{latin}

سپس، یک نمونه جدید از پوسته را فراخوانی می‌کنیم:

\begin{latin}
\begin{lstlisting}
[me@linuxbox ~]$ bash
[me@linuxbox ~]$
\end{lstlisting}
\end{latin}

در ظاهر تفاوت محسوسی مشاهده نمی‌شود، اما در واقعیت یک زیرپوسته (\en{Subshell}) یا نمونه جدید از پوسته اجرا شده است. با بهره‌گیری از دستور \cmd{ps} می‌توان این موضوع را اثبات کرد:

\begin{latin}
\begin{lstlisting}
[me@linuxbox ~]$ ps
    PID TTY          TIME CMD
1011638 pts/9    00:00:00 bash
1011650 pts/9    00:00:00 bash
1011662 pts/9    00:00:00 ps
\end{lstlisting}
\end{latin}

خروجی فوق نشان‌دهنده اجرای هم‌زمان دو نسخه از \cmd{bash} است. از آنجا که نسخه دوم در پس‌زمینه (\en{Background}) قرار داده نشده، هم‌اکنون به‌عنوان یک فرآیند پیش‌زمینه (\en{Foreground Task}) در حال اجراست و پوسته اصلی تا زمان خاتمه یافتن این فرآیند فرزند، در حالت انتظار باقی می‌ماند.

اکنون تلاش می‌کنیم مقدار متغیر \cmd{foo} را که در مرحله قبل مقداردهی کرده بودیم، فراخوانی کنیم:

\begin{latin}
\begin{lstlisting}
[me@linuxbox ~]$ echo $foo

[me@linuxbox ~]$
\end{lstlisting}
\end{latin}

خروجی تهی است؛ زیرا متغیر \cmd{foo} در این نمونه جدید از پوسته تعریف نشده است. به‌طور کلی، متغیرهای پوسته (\en{Shell Variables}) هنگام ایجاد یک فرآیند فرزند، کپی یا منتقل نمی‌شوند؛ در حالی که «متغیرهای محیطی» (\en{Environment Variables}) به‌طور کامل به محیطِ فرآیند فرزند انتقال می‌یابند.

برای بررسی یک ویژگی کلیدی دیگر در رابطه با فرآیندها و محیط آن‌ها، متغیر \cmd{foo} را در همین نشستِ فعلی (فرزند) تعریف می‌کنیم:

\begin{latin}
\begin{lstlisting}
[me@linuxbox ~]$ foo="barbar"
\end{lstlisting}
\end{latin}

سپس با استفاده از دستور \cmd{exit} از این نسخه از \cmd{bash} خارج شده و به نمونه اصلی باز می‌گردیم؛ همان پوسته‌ای که تا کنون در انتظار پایان یافتن فرآیند فرزند خود مانده بود:

\begin{latin}
\begin{lstlisting}
[me@linuxbox ~]$ exit
[me@linuxbox ~]$
\end{lstlisting}
\end{latin}

دوباره دستور \cmd{ps} را اجرا می‌کنیم تا تأیید شود که به نمونه اولیه بازگشته‌ایم:

\begin{latin}
\begin{lstlisting}
[me@linuxbox ~]$ ps
    PID TTY          TIME CMD
1011638 pts/9    00:00:00 bash
1014900 pts/9    00:00:00 ps
\end{lstlisting}
\end{latin}

اکنون مقدار متغیر \cmd{foo} را بررسی می‌کنیم:

\begin{latin}
\begin{lstlisting}
[me@linuxbox ~]$ echo $foo
bar
[me@linuxbox ~]$
\end{lstlisting}
\end{latin}

مشاهده می‌شود که مقدار اولیه (تنظیم شده در نمونه اول) همچنان پابرجاست و تغییرات اعمال شده در فرآیند فرزند، هیچ اثری بر آن نداشته است. این آزمایش گویای یک اصل بنیادین در سیستم‌عامل است: یک فرآیند فرزند هرگز نمی‌تواند محیط فرآیند والد خود را تغییر دهد. با خاتمه یافتن فرآیند فرزند، محیط اختصاصی آن نیز از بین می‌رود. درک این موضوع در هنگام توسعه اسکریپت‌های پوسته (\en{Shell Scripts}) اهمیت حیاتی خواهد داشت.

\section{اجرای برنامه با محیط موقت}

پوسته قابلیت کاربردی دیگری را نیز فراهم می‌کند: توانایی اجرای یک دستور با تخصیص یک متغیر محیطی موقت. در بسیاری از موارد، نیاز است برنامه‌ای را تنها برای یک مرتبه با مقداری خاص در محیط (\en{Environment}) آن اجرا کنیم، بدون آنکه متغیرهای نشست فعلی تغییر کنند.

دستور \cmd{man} مثال بارزی برای این سناریو است. این ابزار به دنبال متغیر محیطی خاصی به نام \cmd{MANWIDTH} می‌گردد که تعیین‌کننده عرض (تعداد کاراکترها) در قالب‌بندی خروجی است.

برای نمونه، جهت نمایش صفحه راهنمای دستور \cmd{ls} با عرض محدود به ۷۵ کاراکتر (به‌منظور ارتقای خوانایی)، می‌توان به شیوه زیر عمل کرد:

\begin{latin}
\begin{lstlisting}
[me@linuxbox ~]$ MANWIDTH=75 man ls
\end{lstlisting}
\end{latin}

این دستور، صفحه راهنمای \cmd{ls} را با عرض بهینه و ساختاری منسجم نمایش می‌دهد.

شایان ذکر است که این قابلیت، فرصتی عالی برای تعریف یک نام مستعار (\en{Alias}) دائمی فراهم می‌کند:

\begin{latin}
\begin{lstlisting}
[me@linuxbox ~]$ alias man='MANWIDTH=75 man'
\end{lstlisting}
\end{latin}

با ثبت این دستور در خط فرمان، از این پس تا پایان نشست جاری، هر بار که فرمان \cmd{man} فراخوانی شود، خروجی آن به‌صورت خودکار بر روی عرض ۷۵ کاراکتر تنظیم خواهد شد.

برای دائمی کردن این تنظیم، باید دستور \cmd{alias} را به فایل \file{.bashrc} اضافه کنید. این فایل، فایل پیکربندی اصلی برای نشست‌های غیرورود تعاملی (مانند باز کردن ترمینال در محیط گرافیکی) است. اگرچه \file{.bashrc} به‌طور پیش‌فرض در نشست‌های ورود (مانند لاگین از طریق کنسول یا \en{SSH}) اجرا نمی‌شود، اما در بسیاری از توزیع‌ها، فایل \file{.bash\_profile} به‌گونه‌ای پیکربندی شده است که \file{.bashrc} را فراخوانی (\en{Source}) کند. با این کار، این نام مستعار در تمام نشست‌های بعدی ترمینال (چه ورود و چه غیرورود) در دسترس خواهد بود.

\section{پیکربندی و شخصی‌سازی محیط کاربری}

اکنون که با مکان و محتوای فایل‌های آغازین آشنا شده‌ایم، می‌توانیم جهت سفارشی‌سازی محیط عملیاتی خود، در آن‌ها تغییرات لازم را اعمال کنیم.

به‌طور کلی، برای افزودن مسیرهای جدید به متغیر \cmd{PATH} یا تعریف متغیرهای محیطی (\en{Environment Variables}) جدید، تغییرات را در فایل \file{.bash\_profile} اعمال کنید (یا معادل آن در توزیع مورد استفاده؛ برای مثال در اوبونتو از فایل \file{.profile} استفاده می‌شود). برای سایر شخصی‌سازی‌ها، از جمله تعریف نام‌های مستعار (\en{Aliases}) یا توابع پوسته، فایل \file{.bashrc} گزینه مناسب‌تری است.

\begin{note}
\textbf{نکته:} همواره توصیه می‌شود تغییرات خود را به فایل‌های موجود در دایرکتوری خانگی (\cmd{\$HOME}) محدود کنید، مگر آنکه به عنوان مدیر سیستم (\en{SysAdmin}) قصد تغییر تنظیمات پیش‌فرض برای تمامی کاربران را داشته باشید. اگرچه ویرایش فایل‌های سراسری مانند \file{/etc/profile} در موارد خاص توجیه فنی دارد، اما جهت حفظ پایداری و امنیت سیستم، در گام‌های نخست از تغییر آن‌ها صرف‌نظر می‌کنیم.
\end{note}

\subsection{ویرایشگرهای متن}

برای ویرایش فایل‌های آغازین پوسته و اغلب فایل‌های پیکربندی سیستم، از برنامه‌ای به نام «ویرایشگر متن» استفاده می‌شود. ویرایشگرهای متن از برخی جهات با واژه‌پردازها (\en{Word Processors}) شباهت دارند؛ هر دو امکان ویرایش کلمات روی صفحه را از طریق مکان‌نمای متحرک فراهم می‌کنند. با این حال، تفاوت بنیادین آن‌ها در این است که ویرایشگرهای متن منحصراً با «متن ساده» (\en{Plain Text}) سروکار دارند و فاقد قالب‌بندی‌های بصری پیچیده هستند. در عوض، ابزارهایی تخصصی برای کدنویسی، پیکربندی سیستم و مدیریت محیط‌های اجرایی در اختیار کاربر قرار می‌دهند. به همین دلیل، ویرایشگرهای متن به‌عنوان ابزاری محوری برای توسعه‌دهندگان نرم‌افزار و مدیران سیستم (\en{SysAdmins}) در محیط لینوکس شناخته می‌شوند و نقش کلیدی در کنترل و تنظیمات سیستم ایفا می‌کنند.

تنوع گسترده‌ای از ویرایشگرهای متن در دنیای لینوکس وجود دارد و اغلب سیستم‌ها به‌صورت پیش‌فرض چندین مورد از آن‌ها را نصب دارند. علت این تکثر، علاقه برنامه‌نویسان به بهینه‌سازی ابزار کار شخصی خود و بازتاب علایق فنی‌شان در عملکرد این نرم‌افزارهاست.

ویرایشگرهای متن به دو دسته‌بندی کلی تقسیم می‌شوند: گرافیکی (\en{GUI}) و مبتنی بر متن (\en{CLI}). محیط‌های دسکتاپ مانند \en{GNOME} و \en{KDE} ویرایشگرهای گرافیکی محبوبی را ارائه می‌دهند؛ \en{GNOME} به‌صورت استاندارد با \cmd{gedit} (که در منوها با عنوان \en{``Text Editor''} شناخته می‌شود) عرضه می‌گردد، در حالی که \en{KDE} دارای دو ویرایشگر اصلی است: \cmd{KWrite} که ویرایشگری ساده و سبک برای کارهای روزمره محسوب می‌شود، و \cmd{Kate} (مخفف \en{KDE Advanced Text Editor}) که امکانات پیشرفته‌تری را برای توسعه‌دهندگان و کاربران حرفه‌ای فراهم می‌آورد.

در مقابل، ویرایشگرهای مبتنی بر متن (ترمینالی) قرار دارند که در محیط‌های فاقد رابط گرافیکی یا ارتباطات راه دور (\en{SSH}) حیاتی هستند. پرکاربردترینِ این ابزارها عبارتند از:

\begin{itemize}
  \item \cmd{nano}: ویرایشگری ساده و کاربرپسند که به‌عنوان جایگزینِ ابزار قدیمی \cmd{pico} طراحی شده است و برای کاربران تازه‌وارد انتخاب ایده‌آلی محسوب می‌شود.
  \item \cmd{vi}: ویرایشگر کلاسیک و سنتی سیستم‌های شبه‌یونیکس (\en{Unix-like}) که در اکثر توزیع‌های مدرن با نسخه پیشرفته‌تری به نام \cmd{vim} (مخفف \en{vi improved}) جایگزین شده است. این ابزار قدرتمند در فصل بعدی به‌طور مفصل بررسی خواهد شد.
  \item \cmd{emacs}: ابزاری عظیم، همه‌کاره و یک محیط جامع برای برنامه‌نویسی که توسط ریچارد استالمن (بنیان‌گذار بنیاد نرم‌افزار آزاد) خلق شده است. با وجود قدرت بی‌نظیر، این برنامه معمولاً به‌صورت پیش‌فرض در اکثر توزیع‌ها نصب نیست.
\end{itemize}

\subsection{کار با ویرایشگرهای متن}

نحوه فراخوانی ویرایشگرهای متن در خط فرمان بدین صورت است که نام ویرایشگر به همراه نام فایل مورد نظر به عنوان آرگومان وارد می‌شود. چنانچه فایل در مسیر مشخص‌شده وجود نداشته باشد، ویرایشگر به‌صورت خودکار فایلی جدید با همان نام ایجاد می‌کند. به عنوان نمونه، برای استفاده از ویرایشگر گرافیکی \cmd{gedit} دستور زیر را اجرا می‌کنیم:

\begin{latin}
\begin{lstlisting}
[me@linuxbox ~]$ gedit some_file
\end{lstlisting}
\end{latin}

این فرمان، رابط گرافیکی \cmd{gedit} را باز کرده و در صورت وجود فایلی با نام \file{some\_file} در دایرکتوری جاری، محتوای آن را جهت ویرایش بارگذاری می‌کند.

از آنجا که سازوکار ویرایشگرهای گرافیکی تا حد زیادی بصری و بی‌نیاز از توضیح است، از بررسی آن‌ها صرف‌نظر کرده و بر نخستین ویرایشگر متنیِ مبتنی بر ترمینال، یعنی \cmd{nano} تمرکز می‌کنیم. پیش از اجرای \cmd{nano} و ویرایش فایل \file{.bashrc}، لازم است اصل رایانش امن (\en{Secure Computing}) را رعایت کنیم؛ بدین معنا که پیش از اعمال هرگونه تغییر در فایل‌های پیکربندی حیاتی، همواره یک نسخه پشتیبان (\en{Backup}) از آن‌ها تهیه کنیم تا در صورت بروز خطای احتمالی، امکان بازگردانی تنظیمات فراهم باشد.

برای ایجاد نسخه پشتیبان از فایل \file{.bashrc}، دستور زیر را اجرا می‌کنیم:

\begin{latin}
\begin{lstlisting}
[me@linuxbox ~]$ cp .bashrc .bashrc.bak
\end{lstlisting}
\end{latin}

انتخاب نام برای نسخه پشتیبان اختیاری است، اما استفاده از پسوندهای استانداردی نظیر \file{.bak}، \file{.sav}، \file{.old} یا \file{.orig} جهت تشخیص فایل‌های پشتیبان در ادبیات لینوکس بسیار رایج است. همچنین به یاد داشته باشید که دستور \cmd{cp} به‌صورت پیش‌فرض و بدون هشدار قبلی، فایل‌های هم‌نام در مقصد را بازنویسی (\en{Overwrite}) می‌کند.

اکنون که نسخه پشتیبان را ایجاد کرده‌ایم، ویرایشگر را اجرا می‌کنیم:

\begin{latin}
\begin{lstlisting}
[me@linuxbox ~]$ nano .bashrc
\end{lstlisting}
\end{latin}

پس از اجرای \cmd{nano} با محیطی مشابه تصویر زیر مواجه خواهید شد:

\begin{latin}
\begin{lstlisting}
GNU nano 2.0.3              File: .bashrc                    .  .     .

# .bashrc

# Source global definitions
if [ -f /etc/bashrc ]; then
    . /etc/bashrc
fi

# User specific aliases and functions


                         [ Read 8 lines ]
^G Help        ^O Write Out   ^F Where Is    ^K Cut         ^T Execute     ^C Location    M-U Undo
^X Exit        ^R Read File   ^\ Replace     ^U Paste       ^J Justify     ^/ Go To Line  M-E Redo
\end{lstlisting}
\end{latin}

\begin{note}
\textbf{نکته:} چنانچه ویرایشگر \cmd{nano} بر روی سیستم شما نصب نیست، می‌توانید به‌طور موقت از یک ویرایشگر گرافیکی استفاده کنید.
\end{note}

ساختار رابط کاربری در این برنامه شامل سه بخش اصلی است: سرآیند (\en{Header}) در بالاترین قسمت، بدنه محتوای متنی فایل در مرکز، و فهرست فرامین (\en{Command Menu}) در پایین صفحه. از آنجا که \cmd{nano} در اصل به‌عنوان جایگزینی برای ویرایشگر ساده‌ی یک کلاینت ایمیل طراحی شده است، ابزارهای ویرایشی آن بسیار مینیمال و فاقد پیچیدگی‌های رایج در ویرایشگرهای پیشرفته است.

نخستین و حیاتی‌ترین مهارتی که باید در مواجهه با هر ویرایشگر متن بیاموزید، چگونگی خروج از محیط برنامه است. در \cmd{nano}، جهت خروج از کلید ترکیبی \cmd{Ctrl+x} استفاده می‌شود. این دستور در منوی راهنمای پایین صفحه با نماد \cmd{\^{}X} درج شده است؛ در ادبیات سیستم‌های کامپیوتری، نویسه‌ی \cmd{\^{}} (\en{Caret}) نمایانگر کلید \en{Ctrl} است و این نوع نمادگذاری برای «نویسه‌های کنترلی» (\en{Control Characters}) در بسیاری از محیط‌های نرم‌افزاری استاندارد محسوب می‌شود.

دومین دستور کلیدی که باید به خاطر بسپارید، ذخیره‌سازی تغییرات (\en{Save}) است که در \cmd{nano} با کلید ترکیبی \cmd{Ctrl+o} انجام می‌شود. با تسلط بر این دو فرمان، آمادگی لازم برای ویرایش فایل را دارید. با استفاده از کلیدهای جهت‌نما یا \en{PageDown}، مکان‌نما را به انتهای فایل هدایت کرده و سطر‌های زیر را به فایل \file{.bashrc} اضافه کنید:

\begin{latin}
\begin{lstlisting}
umask 0002
export HISTCONTROL=ignoredups
export HISTSIZE=1000
alias l='ls -d .* --color=auto'
alias ll='ls -l --color=auto'
\end{lstlisting}
\end{latin}

\begin{note}
\textbf{نکته:} احتمال دارد برخی از این خطوط به‌صورت پیش‌فرض در توزیع لینوکس شما موجود باشد؛ با این حال، تکرار آن‌ها در انتهای فایل تداخلی در عملکرد سیستم ایجاد نخواهد کرد.
\end{note}

در جدول زیر، شرح فنی و عملکردی هر یک از دستورهای افزوده‌شده تبیین شده است:

\begin{table}[htbp]
\centering
\caption{پیکربندی‌های افزوده‌شده به فایل \file{.bashrc}}
\label{tab:bashrc-configs}
\begin{tabularx}{\textwidth}{l X}
\toprule
\textbf{دستور} & \textbf{معنا} \\
\midrule
\cmd{umask 0002} & تنظیم ماسک ایجاد فایل (\cmd{umask}) جهت رفع محدودیت‌های دسترسی در دایرکتوری‌های اشتراکی. این مقدار، مجوزهای پیش‌فرض را به‌گونه‌ای تنظیم می‌کند که فایل‌های جدید برای اعضای گروه قابل نوشتن باشند (مطابق با مباحث فصل ۹). \\
\addlinespace
\cmd{export HISTCONTROL=ignoredups} & تنظیم متغیر محیطی \cmd{HISTCONTROL} با مقدار \cmd{ignoredups} که مانع از ذخیره‌سازی دستورات تکراری و متوالی در فایل تاریخچه می‌شود. \\
\addlinespace
\cmd{export HISTSIZE=1000} & ارتقای ظرفیت ذخیره‌سازی تاریخچه دستورها به ۱۰۰۰ سطر. \\
\addlinespace
\cmd{alias l.='ls -d .* --color=auto'} & تعریف نام مستعار \cmd{l.} جهت نمایش فایل‌ها و دایرکتوری‌های پنهان (مواردی که با نقطه شروع می‌شوند)، همراه با دایرکتوری جاری (\cmd{.}) و دایرکتوری والد (\cmd{..}). \\
\addlinespace
\cmd{alias ll='ls -l --color=auto'} & تعریف نام مستعار \cmd{ll} برای مشاهده فهرست محتویات دایرکتوری با جزئیات کامل (\en{Long Listing Format}). \\
\bottomrule
\end{tabularx}
\end{table}

همان‌طور که مشاهده می‌شود، بسیاری از این دستورات فاقد وضوح معنایی برای کاربر هستند؛ لذا جهت ارتقای خوانایی اسکریپت، شایسته است توضیحاتی (\en{Comments}) را به فایل \file{.bashrc} بیافزاییم. با استفاده از ویرایشگر، قطعه‌کد قبلی را به صورت زیر اصلاح نمایید:

\begin{latin}
\begin{lstlisting}
# Change umask to make directory sharing easier
umask 0002

# Ignore duplicates in command history and increase
# history size to 1000 lines
export HISTCONTROL=ignoredups
export HISTSIZE=1000

# Add some helpful aliases
alias l.='ls -d .* --color=auto'
alias ll='ls -l --color=auto'
\end{lstlisting}
\end{latin}

اکنون ساختار فایل بسیار گویاتر شده است. پس از اعمال این اصلاحات، با فشردن کلیدهای \cmd{Ctrl+o} تغییرات را در فایل \file{.bashrc} ذخیره نموده و سپس با استفاده از \cmd{Ctrl+x} از محیط \cmd{nano} خارج شوید.

\subsection{مستندسازی و درج توضیحات}

هنگام ویرایش فایل‌های پیکربندی، همواره شایسته است تغییرات اعمال‌شده را از طریق درج توضیحات (\en{Comments}) مستندسازی کنید. شاید جزئیات این تغییرات در لحظه برای شما شفاف باشد، اما بازخوانی آن‌ها پس از گذشت چندین ماه بدون وجود راهنما دشوار خواهد بود. بنابراین، با نگارش توضیحات دقیق، فرآیند عیب‌یابی و مدیریت سیستم را در آینده برای خود تسهیل کنید. همچنین، ثبت تاریخچه‌ای از تغییرات سیستمی در قالب یک گزارش مجزا، رویکردی حرفه‌ای در مدیریت سیستم محسوب می‌شود.

در اسکریپت‌های پوسته و فایل‌های آغازین \cmd{bash} (مانند \file{\textasciitilde/.bashrc})، از نماد \cmd{\#} برای شروع یک سطر توضیحی استفاده می‌شود. شایان ذکر است که سایر فایل‌های پیکربندی ممکن است از کاراکترهای متفاوتی برای این منظور بهره ببرند. اغلب این فایل‌ها به‌صورت پیش‌فرض حاوی توضیحات راهنما هستند؛ توصیه می‌شود از این متون به‌عنوان الگویی برای درک ساختار پیکربندی استفاده کنید.

در بسیاری از فایل‌های پیکربندی، خطوطی وجود دارند که با هدف جلوگیری از اجرا توسط برنامه، «کامنت» شده‌اند. این تکنیک عمدتاً برای ارائه پیشنهادهایی جهت تنظیمات احتمالی یا نمایش سینتکس (نحو) صحیح دستورها به کار می‌رود. برای نمونه، در فایل \file{.bashrc} توزیع اوبونتو ۱۸.۰۴، سطر‌های زیر مشاهده می‌شود:

\begin{latin}
\begin{lstlisting}
# some more ls aliases
#alias ll='ls -l'
#alias la='ls -A'
#alias l='ls -CF'
\end{lstlisting}
\end{latin}

سه سطر پایانی، دستورهای معتبر جهت تعریف نام مستعار (\en{Alias}) هستند که به دلیل وجود کاراکتر \cmd{\#} در ابتدای آن‌ها، غیرفعال شده‌اند. با حذف این علامت (فرآیندی که در ادبیات فنی \en{Uncommenting} نامیده می‌شود)، این نام‌های مستعار فعال خواهند شد. به‌همین ترتیب، با افزودن نماد \cmd{\#} به ابتدای هر خط پیکربندی، می‌توانید آن را بدون حذفِ دائمی اطلاعات، به‌صورت موقت از چرخه اجرا خارج کنید.

\subsection{اعمال و فعال‌سازی تغییرات}

تغییراتی که در فایل \file{.bashrc} اعمال کردیم، به‌صورت خودکار در نشست (\en{Session}) فعلی ترمینال فعال نخواهند شد؛ چرا که این فایل تنها در لحظه‌ی آغازینِ بارگذاری پوسته بازخوانی می‌شود. برای اعمال این اصلاحات بدون نیاز به بستن ترمینال و ورود مجدد، می‌توان پوسته \cmd{bash} را وادار به بازخوانی و اجرای آنیِ فایل پیکربندی نمود. این عمل با استفاده از دستور \cmd{source} به شرح زیر انجام می‌پذیرد:

\begin{latin}
\begin{lstlisting}
[me@linuxbox ~]$ source ~/.bashrc
\end{lstlisting}
\end{latin}

پس از اجرای فرمان فوق، پیکربندی‌های جدید در نشست جاری تزریق می‌شوند. اکنون می‌توانید صحت عملکرد یکی از نام‌های مستعار (\en{Alias}) تعریف‌شده را بررسی کنید:

\begin{latin}
\begin{lstlisting}
[me@linuxbox ~]$ ll
\end{lstlisting}
\end{latin}

با اجرای این دستور، باید فهرست فایل‌ها را با جزئیات کامل (\en{Long Format}) مشاهده کنید که نشان‌دهنده موفقیت‌آمیز بودن فرآیند است.

\subsection{دستور source و ماهیت فایل‌های آغازین}

دستور \cmd{source} (که با نماد نقطه \cmd{.} نیز شناخته می‌شود) یکی از فرامین داخلی (\en{Built-in}) پوسته است. وظیفه‌ی این دستور، خواندن و اجرای مستقیم محتویات یک فایل در محیط نشست فعلی است؛ به‌گونه‌ای که گویی تمامی خطوط آن فایل، همان لحظه به‌صورت دستی توسط کاربر در خط فرمان تایپ شده باشند. این موضوع بیانگر نکته‌ی مهمی است: تمامی آن ساختارهای به‌ظاهر پیچیده که در فایل‌های آغازین (\en{Startup Files}) مشاهده کردیم، در واقع مجموعه‌ای از دستورهای استاندارد هستند که پوسته به‌طور بومی قادر به تفسیر و اجرای آن‌هاست.

بسیاری از سیستم‌عامل‌های قدیمی مبتنی بر متن (مانند \en{DOS} یا \en{CP/M})، محیط‌هایی محدود بودند که عمدتاً وظیفه‌ی ساده‌ی فراخوانی و اجرای برنامه‌ها را بر عهده داشتند. اگرچه پوسته‌های یونیکسی نیز به بهترین شکل این وظیفه را انجام می‌دهند، اما فراتر از یک واسط اجرای برنامه عمل می‌کنند. این پوسته‌ها در حقیقت محیط‌هایی بسیار غنی و پویا هستند که با بهره‌گیری از متغیرهای محیطی، نام‌های مستعار و توابع، امکان شخصی‌سازی بی‌نظیری را برای کاربر فراهم می‌آورند.

\section{جمع‌بندی}

در این فصل، یکی از مهارت‌های بنیادین در دنیای لینوکس، یعنی ویرایش فایل‌های پیکربندی از طریق ویرایشگرهای متن را فراگرفتیم. از این پس، هنگام مطالعه صفحات راهنما (\en{man pages})، به متغیرهای محیطی (\en{Environment Variables}) که توسط دستورات مختلف پشتیبانی می‌شوند، دقت ویژه‌ای داشته باشید؛ چرا که ابزارهای قدرتمندی برای شخصی‌سازی رفتار برنامه‌ها در اختیار شما قرار می‌دهند. در فصول آتی، با مفاهیم «توابع پوسته» (\en{Shell Functions}) آشنا خواهیم شد؛ قابلیتی بسیار منعطف که به شما اجازه می‌دهد مجموعه‌ای از فرامین سفارشی و پیچیده را در فایل‌های آغازین \cmd{bash} تعریف و اجرا کنید.

\section{منابع مطالعه تکمیلی}

برای درک عمیق و جزئی رفتار پوسته هنگام آغاز به کار (\en{startup})، مطالعه بخش \en{INVOCATION} در صفحه راهنمای \en{Bash} توصیه می‌شود. این بخش، ترتیب و شرایط بارگذاری فایل‌های پیکربندی را در انواع نشست‌های پوسته به‌دقت تشریح می‌کند.

\begin{latin}
\begin{lstlisting}
[me@linuxbox ~]$ man bash
\end{lstlisting}
\end{latin}